\section{Introduction}
\label{sec:intro}

A researcher who tries many variants of a trading strategy and reports the best one reports a biased number. The best of $N$ strategies with no skill has an in-sample Sharpe ratio that grows with $N$, roughly in proportion to $\sqrt{2\ln N}$ times the spread of the individual estimates \citep{bailey2014pseudo}, and the same selection effect pervades published factor research \citep{harvey2016cross}. Several corrections are now standard: the deflated Sharpe ratio \citep{bailey2014dsr}, built on the probabilistic Sharpe ratio \citep{bailey2012frontier}; the multiple-testing haircut of \citet{harvey2015backtesting}; family-wise procedures such as those of \citet{holm1979} and \citet{sidak1967}; and resampling tests in the tradition of \citet{white2000reality}, \citet{hansen2005spa} and \citet{romano2005stepwise}.

Each correction rests on assumptions---normal or independent returns, independent trials, a known number of trials---that real searches violate in predictable ways. Returns are fat-tailed, skewed, volatility-clustered and sometimes autocorrelated; strategy variants share most of their positions and so are correlated. Inference for a single Sharpe ratio under such violations is well studied \citep{lo2002sharpe,opdyke2007comparing,ledoit2008robust}. How much each violation distorts each \emph{selection} correction, applied to the best of a search, is rarely answered in a form a practitioner can look up.

This paper answers it with a simulation benchmark, the \emph{Null Zoo}. The benchmark draws complete research searches---\nzTrials\ candidate strategies of 504 daily returns each---from eight return families whose ground truth is known by construction, applies seven validators as their implementations are meant to be called, and records how often each one declares the best strategy skilled. When no strategy has skill, that frequency is the validator's \emph{size}, which should equal its nominal level. When one strategy has a true annualized Sharpe ratio of 2, it is the validator's \emph{power}, which we report both raw and size-adjusted. We follow the ADEMP structure of \citet{morris2019simulation} for planning and reporting the study; every rate's Monte Carlo standard error follows from the stated number of replications.

Four findings stand out.
\begin{enumerate}
  \item Positive autocorrelation breaks the tests that ignore it. With a lag-one autocorrelation of 0.2, the \v{S}id\'ak-adjusted $t$ test, the haircut, the non-normal-error test and an i.i.d.\ bootstrap all rejected between \nzBootstrapBestArOneSizeExt\% and \nzLuckTrialsArOneSizeExt\% of skill-less searches at a nominal 5\%. Deflating the $t$-statistic by Lo's autocorrelation-adjusted annualisation factor brings the size to \nzLuckTrialsLoArOneSizeExt\%, in the case where the factor's assumption holds exactly, at a cost of a slight excess size elsewhere.
  \item Under negative skew the Student-$t$ tests reject \nzLuckTrialsSkewNegativeSizeExt\% at a nominal 5\%. The non-normal standard error and the bootstrap each fix one shape and break another. Under fat tails, selection favours series with a large positive outlier (the selected series' mean skewness is \nzSkewBestStudentTFour, against \nzSkewOthersStudentTFour\ for the rest), which the non-normal standard error then mistakes for a property of the returns. No validator is calibrated across skew and fat tails.
  \item The deflated Sharpe ratio, with the 0.95 acceptance rule its authors use, rejects at most \nzDeflatedSharpeCorrelatedTrialsSizeExt\% of skill-less searches and, under i.i.d.\ normal returns, fails to accept \nzDsrFailToAccept\% of searches that contain a strategy with a true Sharpe ratio of 2 over two years. The rule tests a strictly harder hypothesis than ``the best strategy has no skill'', and the skilled strategy raises its own benchmark. Size-adjusted, it recovers most of the lost power under i.i.d.\ returns (\nzDeflatedSharpeIidNormalAdjPowerExt\% against \nzLuckTrialsIidNormalAdjPowerExt\%) but trails the calibrated test by \nzDsrAdjGapMin--\nzDsrAdjGapMax\ points across families: the threshold costs most of the power, the statistic the rest.
  \item A bootstrap's number of resamples materially affects its power. With 400 resamples and a \v{S}id\'ak adjustment for 20 trials, the test can reject only when no resample reaches the observed Sharpe ratio. On the same searches, raising the resamples to \nzBootDraws\ raises its power from \nzBrIidNormalPowerBootFourHundred\% to \nzBrIidNormalPowerBootMax\% under i.i.d.\ normal returns, with almost no change in size.
\end{enumerate}
We also note an analytic point that affects how such comparisons are read: the haircut of \citet{harvey2015backtesting} is defined with two-sided $p$-values, and read one-sided it is almost exactly the \v{S}id\'ak-adjusted $t$ test, so comparisons at unmatched sides measure the choice of sides.

The benchmark scores a statistic introduced by its own author, the \emph{luck-equivalent trials} count. As a test it is the \v{S}id\'ak-adjusted Student $t$ test; it reports the same evidence in units of trials. \citet{bartos2025living} argue that simulation studies designed by a method's developers carry a conflict of interest. We report that conflict, publish every generator, validator and seed, and note that the benchmark's clearest failure---autocorrelation---applies to the author's statistic as much as to the others. The published run reproduces byte for byte from its code, and a robustness run with a different generator, five times as many replications and a finer bootstrap shows that the findings are not artefacts of the generator or of Monte Carlo error.

\paragraph{Related work.}
\citet{arian2024backtest} compare out-of-sample testing schemes---walk-forward, $k$-fold, purged and combinatorial purged cross-validation---in a synthetic market with known structure and score them by the probability of backtest overfitting \citep{bailey2017pbo} and the deflated Sharpe ratio. Their object is the resampling scheme used to produce a backtest; ours is the statistical correction applied to its result. \citet{santoni2026minerva} combine several of the same corrections into one robustness grade and report its discrimination on synthetic markets, including a baseline for the deflated Sharpe ratio alone; they do not report size at a nominal level or power for the individual corrections, which is what a user choosing among them needs. The size and power of multiple-testing procedures under dependence are well studied in statistics; our contribution is a small, open, versioned benchmark that uses the return shapes and search structures of trading research and calls each correction the way it is called in practice.
